\begin{statement}[C.6.3; Proved: the CM models and their height identities]
The eleven rational CM fibres have the global minimal models,
section coordinates, component classes, and finite local heights
in the following tables. On each model $a_1=a_2=0$, and we
write $[a_3,a_4,a_6]$ for
$y^2+a_3y=x^3+a_4x+a_6$. The displayed point is $(0,y_P)$.
The sign of $P$, which is immaterial to all these heights,
is chosen by the recorded change of variables.
The constants $A_t$ in increasing row order are
\[
 \sqrt6,\quad 2^{2/3}3^{1/4},\quad 2^{3/4},\quad
 2,\quad 2,\quad \sqrt3,\quad 1,\quad 3^{5/12},\quad
 1,\quad1,\quad1.
\]
\end{statement}
\begin{proof}
We substitute the eleven exact values $t=j/1728$, and perform
the integral minimal-model transformations and Tate reduction
algorithm. The transformations and resulting invariants are
printed in \path{code/C_cm_models.out}; the exact valuation
computations are in \path{code/C_cm_local_heights.out}.
For each bad prime the program computes $mP$ by the rational
group law, tests the two partial derivatives of the minimal
cubic, and evaluates $\psi_m$ by its exact recursion. Formula
\eqref{eq:ch-local-exact} then proves each entry. Outside the
listed primes the discriminant is a unit and $P$ is integral,
so the local height is zero. Multiplying $p^{\lambda_p/\log p}$
gives the listed constants $A_t$.

We specify the component labels without relying on an
unrecorded ordering of a resolution graph. The label $0$
means the identity component, and $c_P$ is the class represented
by the explicitly displayed section $P$. At an odd prime of
type $I_0^*$ this nonidentity class is the outer component
indexed by the root zero of the residual Tate cubic; the
scaling $x=pX$ identifies that root. When there are two
rational components, $c_P$ is their unique nonzero class.
The only remaining ambiguity is row 5 at $2$, where the
component group is $(\Z/2)^2$. Here let $T=(-14,0)$;
we label it with the basis $(c_P,c_T)$. Indeed
$P+T=(18,64)$ after the corresponding choice of sign,
and $P,T,P+T$ all reduce to the singular point, whereas
$2P=(25/4,99/8)$ reduces to the identity component.
Thus the three nonzero classes are explicitly separated.

The independent analytic check uses the exact quadratic CM
period ratio, solves $\wp(\mu)=0$ by its convergent Fourier
series, and evaluates the theta product at 270-digit working
precision. The value of $j$ agrees with its prescribed integer
to absolute error less than $10^{-245}$. For every row the
direct value of $\log N$ agrees with
$\log A_t-\widehat h_O(P_t)$ to more than 245 decimal places.
These numerical comparisons test the exact normalization;
the algebraic local-height entries were already proved by
integer and rational arithmetic. The full precision values
are in \path{records/C_cm_theta_pslq.json}.
\end{proof}
